\section[Bell Inequalities and tau-n]{Bell Inequalities and \texorpdfstring{$\tau_n$}{tau_n}}
\begin{quote}\small
\noindent\textbf{Status.} \emph{Legacy physical-readout block inside the active main body.}\par
\noindent\textbf{Block function.} This section places the Bell/nonlocality discussion on the same readout-relative footing and makes clear what role the $\tau_n$ layer plays there.
\end{quote}
\label{sec:bell}
%% ============================================================

\begin{theorem}[Bell consistency]
\label{thm:bell-consistency}
The proper time $\tau_n$ is not a local hidden variable in Bell's sense.
The framework is consistent with quantum mechanical violation of Bell
inequalities.
\begin{proof}
Bell's theorem assumes a hidden variable $\lambda$ assigned independently
to each particle: measurement outcomes $A(\lambda_A,a)$ and
$B(\lambda_B,b)$ with $\lambda_A,\lambda_B$ statistically independent.

In our framework $\tau_n$ belongs to the pair $(f,Jf)$, not to each
element separately: $\lambda_A=\lambda_B=\tau_n$.
This is a contextual variable --- Bell's theorem does not exclude
violation by contextual theories.

The correlation function $E(a,b)$ follows from the structure of
$\Hplusplus$: since $C(f)\in\Hplusplus$ carries an irreducible
representation of $\WA$ (Stone--von Neumann, \cite{brendecke2026rh}),
the standard quantum correlations $E(a,b)=\cos\theta_{ab}$ are
reproduced, giving CHSH value $S=2\sqrt{2}>2$.

Hence:
\begin{itemize}
  \item $\tau_n$ explains why the correlation exists: shared proper time,
        fixed at the near-collision;
  \item Stone--von Neumann determines how strong the violation is
        ($S=2\sqrt{2}$).
\end{itemize}
Lorentz invariance and quantum mechanics are consistent.
\end{proof}
\end{theorem}

\begin{remark}[What $\tau_n$ does and does not do]
$\tau_n$ encodes the magnitude $|\zminus|$ of the displacement from the
photon sector at the moment of the near-collision.
It does not determine the measurement direction.
When $f$ is measured, $\zminus$ is projected onto a value; the paired
$Jf$ is projected onto $-\zminus$ (the $J$-image).
The anticorrelation was built in at the collision event, not transmitted
at measurement. No signal travels faster than light.
\end{remark}

%% ============================================================

